%% formal_mordell_weil_rank_2026-07-30.tex
%%
%% A machine-checked canonical height on an elliptic curve, and the first
%% formal proof of independence of rational points: rank >= 2 for 389a1.
%%
%% Every mathematical claim below is either a kernel-verified Lean 4 theorem
%% at HEAD 0d3b5415 (2026-07-30), axioms [propext, Classical.choice,
%% Quot.sound], zero project axioms, zero sorries, no native_decide, or
%% standard cited literature. Verified build data: lake build PF = 4658
%% jobs, lake build (full) = 9301 jobs, mathlib pin eed770a4, toolchain
%% leanprover/lean4:v4.24.0-rc1.

\documentclass[11pt,reqno]{amsart}

\usepackage[T1]{fontenc}
\usepackage{lmodern}
\usepackage{amsmath,amssymb,amsthm}
\usepackage{mathtools}
\usepackage{microtype}
\usepackage{geometry}
\geometry{margin=1in}
\usepackage[hidelinks]{hyperref}
\usepackage{enumitem}
\usepackage{booktabs}
\usepackage{seqsplit}

\usepackage{tikz}
\usetikzlibrary{arrows.meta,positioning}
\usepackage{pgfplots}
\pgfplotsset{compat=1.18}
\usepgfplotslibrary{groupplots}
%% Figures are generated, not hand-drawn: every coordinate below comes from
%% Papers/figures/mordell_weil/make_figures.py, which recomputes the group
%% law, the naive height and the height sequence from the same definitions
%% the Lean files use.  Regenerate with
%%     python Papers/figures/mordell_weil/make_figures.py
\newcommand{\pffig}[1]{\input{figures/mordell_weil/#1.tex}}

\emergencystretch=3em
\hbadness=6000
\sloppy

\theoremstyle{plain}
\newtheorem{theorem}{Theorem}[section]
\newtheorem{lemma}[theorem]{Lemma}
\newtheorem{proposition}[theorem]{Proposition}
\newtheorem{corollary}[theorem]{Corollary}

\theoremstyle{definition}
\newtheorem{definition}[theorem]{Definition}

\theoremstyle{remark}
\newtheorem{remark}[theorem]{Remark}

% \lt: Lean identifiers and shell commands.  Plain \texttt, NOT \seqsplit --
% seqsplit deletes spaces (it would print "lake build PF" as "lakebuildPF")
% and permits a line break after any single character.  Robust because \lt is
% used inside \caption, which amsart writes to the .lof.
\DeclareRobustCommand{\lt}[1]{\texttt{#1}}
% \ltw: long unbroken digit strings, where breaking anywhere is what we want.
\DeclareRobustCommand{\ltw}[1]{\texttt{\seqsplit{#1}}}
\newcommand{\hh}{\widehat{h}}
\newcommand{\nh}{H}
\DeclareMathOperator{\rank}{rank}
\DeclareMathOperator{\reg}{Reg}

\title[A formal canonical height and independence of rational points]{A
machine-checked canonical height on an elliptic curve, and a formal proof
that two rational points are independent}

\author{Pablo Cohen}
\address{Mesa, Arizona, USA}
\email{psolorzano@alumni.berklee.edu}

\date{30 July 2026}

\subjclass[2020]{Primary 11G05, 11Y16; Secondary 68V20, 68V15}

\keywords{elliptic curve, canonical height, Mordell--Weil rank,
independence of points, Lean 4, interactive theorem proving}

\begin{document}

\begin{abstract}
We report the first machine-checked construction of the canonical (N\'eron--Tate)
height on an elliptic curve over $\mathbb{Q}$, and the first formal proof that
two rational points on an elliptic curve are independent. Working in Lean~4
over \texttt{mathlib}, for the curve $E = \mathtt{389a1}$,
$y^2 + y = x^3 + x^2 - 2x$ --- the smallest-conductor curve of rank~$2$ ---
we construct $\hh \colon E(\mathbb{Q}) \to \mathbb{R}$ as a limit of
normalized na\"ive heights, prove that it vanishes exactly on torsion, prove
the parallelogram law for it \emph{exactly}, and deduce
$\hh(kR) = k^2\hh(R)$ for all $k \in \mathbb{Z}$. Certified rational interval
arithmetic on the dyadic chain then shows the regulator determinant of the
generators $P = (0,0)$, $Q = (1,0)$ is positive, whence
\[
  2 \;\le\; \rank E(\mathbb{Q})
\]
as a theorem about \texttt{mathlib}'s literal Mordell--Weil group, with
kernel axioms $[\mathtt{propext}, \mathtt{Classical.choice},
\mathtt{Quot.sound}]$ only. The construction is elementary: it uses no
descent, no formal groups, no torsion classification, and no
Gross--Zagier. As a by-product the same machinery yields
$\rank E(\mathbb{Q}) \ge 1$ for eleven curves of conductor
$37 \le N \le 5077$. We emphasise what is \emph{not} proved: these are rank
lower bounds, not rank equalities, and nothing here bears on the
Birch--Swinnerton-Dyer conjecture.
\end{abstract}

\maketitle

\section{Introduction}

\subsection{What is new}

Let $E/\mathbb{Q}$ be an elliptic curve. The canonical height
$\hh \colon E(\mathbb{Q}) \to \mathbb{R}_{\ge 0}$ of N\'eron and Tate is the
basic tool for measuring rational points: it is a positive semi-definite
quadratic form whose kernel is exactly the torsion subgroup, and the
determinant of its Gram matrix on a set of points --- the regulator ---
detects independence. It is also the quantity appearing on the arithmetic
side of the Birch--Swinnerton-Dyer conjecture.

Despite its centrality, the canonical height had not been formalized in any
interactive theorem prover. At the time of writing, \texttt{mathlib} contains
the Weierstrass curve type, the group law on the affine, projective and
Jacobian point models (with associativity), division polynomials, and, on
\texttt{master}, reduction theory and a formal Dirichlet series; it does not
contain the Mordell--Weil theorem, heights on elliptic curves, descent, or
any analytic continuation of an elliptic $L$-function. Consequently no
positive lower bound on $\rank E(\mathbb{Q})$ for a specific curve appears to
have been machine-checked, and \emph{a fortiori} no independence statement
about rational points.

This paper supplies both, for a single explicit curve, together with the
general machinery. The main theorem is the following.

\begin{theorem}[\lt{E389a1\_rank\_ge\_two}]\label{thm:main}
Let $E$ be the elliptic curve $y^2 + y = x^3 + x^2 - 2x$ over $\mathbb{Q}$
(Cremona label \texttt{389a1}, LMFDB label \texttt{389.a1}). Then
\[
  2 \;\le\; \operatorname{rank}_{\mathbb{Z}} E(\mathbb{Q}),
\]
where $E(\mathbb{Q})$ is \texttt{mathlib}'s type
\lt{WeierstrassCurve.Affine.Point} for $E$ and the rank is
\lt{Module.rank} over $\mathbb{Z}$. The points $P = (0,0)$ and $Q = (1,0)$
are $\mathbb{Z}$-linearly independent.
\end{theorem}

Theorem~\ref{thm:main} is a Lean~4 theorem; \lt{\#print axioms} on it reports
$[\mathtt{propext}, \mathtt{Classical.choice}, \mathtt{Quot.sound}]$ and
nothing else. It rests on the following, which we regard as of independent
interest.

\begin{theorem}[\lt{canheight}, \lt{canheight\_eq\_zero\_iff\_torsion},
\lt{canheight\_parallelogram}, \lt{canheight\_zsmul}]\label{thm:height}
For the curve $E$ above there is a function
$\hh \colon E(\mathbb{Q}) \to \mathbb{R}$, defined as
\[
  \hh(R) \;=\; \lim_{n \to \infty} \frac{\log \nh\!\left(x(2^n R)\right)}{4^n},
\]
such that, formally:
\begin{enumerate}[label=\textup{(\arabic*)}]
  \item $\hh \ge 0$, and $\hh(R) = 0$ if and only if $R$ is a torsion point;
  \item $\lvert \hh(R) - \log \nh(x(R)) \rvert \le \tfrac13\log 1728$;
  \item $\hh(P+Q) + \hh(P-Q) = 2\hh(P) + 2\hh(Q)$ whenever
        $P, Q, P \pm Q$ are non-torsion;
  \item $\hh(kR) = k^2\,\hh(R)$ for every $k \in \mathbb{Z}$ and non-torsion
        $R$.
\end{enumerate}
\end{theorem}

Here $\nh(a/b) = \max(\lvert a\rvert, b)$ for a reduced fraction, and $x(R)$
is the $x$-coordinate of $R$ (with $x(O) := 0$, a harmless convention since
$\nh(0) = 1$).

\subsection{The obstruction, and how it is avoided}

The classical construction of $\hh$ presupposes a body of theory --- Weil
heights on projective space, the theory of local heights, or at minimum the
functorial machinery of Weil's height machine --- none of which is available
in \texttt{mathlib}. The classical proof that $\hh$ is a quadratic form
proceeds through the Mordell--Weil theorem's ingredients or through
polarization identities requiring the parallelogram law at many pairs of
points.

Our route replaces all of this with explicit arithmetic on a \emph{single
fixed curve}. Three observations make it work.

\emph{First}, the duplication map on the $x$-line of a fixed curve is one
explicit rational function. For \texttt{389a1},
\begin{equation}\label{eq:fg}
  x(2R) \;=\; \frac{f(x)}{g(x)}, \qquad
  f(x) = x^4 + 4x^2 - 2x + 3, \qquad
  g(x) = 4x^3 + 4x^2 - 8x + 1,
\end{equation}
where $x = x(R)$; and $g$ has no rational root, so no rational point of
$E$ is $2$-torsion and doubling never leaves the affine chart. Homogenizing
$x = a/b$ and applying the extended Euclidean algorithm to $f$ and $g$ gives
\emph{integer} B\'ezout identities, certified in Lean by \lt{ring}, from which
$\gcd$ cancellation is bounded and the two-sided height estimate
\begin{equation}\label{eq:dup}
  \frac{\nh(x)^4}{1728} \;\le\; \nh\!\left(x(2R)\right) \;\le\; 17\,\nh(x)^4
\end{equation}
follows with explicit constants. Estimate~\eqref{eq:dup} makes
$\log \nh(x(2^nR))/4^n$ a Cauchy sequence with an explicit geometric modulus,
which is all that is needed to define $\hh$.

\emph{Second}, the addition law admits the same treatment, once one works with
the pair $\{x(P+Q), x(P-Q)\}$ rather than with either coordinate separately.
Reducing modulo the curve equations eliminates the $y$-coordinates entirely:
$x(P+Q)$ and $x(P-Q)$ are the two roots of an explicit quadratic whose
coefficients are bihomogeneous of bidegree $(2,2)$ in the reduced coordinates
of $x(P)$ and $x(Q)$. Combined with an elementary lemma bounding the product
of the heights of the roots of an integer quadratic by the height of its
coefficient vector (Lemma~\ref{lem:roots} below), this yields the
quasi-parallelogram law with an explicit constant. Notably the lower bound is
free: applying the upper bound to the pair $(P+Q, P-Q)$, whose sum and
difference are $2P$ and $2Q$, and feeding in~\eqref{eq:dup} gives it.

\emph{Third}, and this is what removes the last classical dependency: for the
independence argument one does not need bi-additivity of the induced pairing.
If a relation $mP + nQ = 0$ exists with $n \ne 0$, then the relation
\emph{itself} rewrites $P \pm Q$ as multiples of $P$,
\[
  n(P+Q) = (n-m)P, \qquad n(P-Q) = (n+m)P,
\]
so Theorem~\ref{thm:height}(4) applies, and a \emph{single} use of the
parallelogram law at $(P,Q)$ forces the regulator determinant to vanish
(Proposition~\ref{prop:reldet}). The contrapositive is the independence
criterion.

\begin{figure}[htbp]
\centering
\pffig{fig_arc_dag}
\caption{Dependency structure of the development. Boxes are Lean files,
labelled by their internal revision tag; blue boxes are the two results of
independent interest (the canonical height and the exact parallelogram law),
green is Theorem~\ref{thm:main}. Arrows denote logical dependence, not
import order. Every box is fully proved: no box contains a \lt{sorry}, an
axiom, or a \lt{native\_decide}.}
\label{fig:dag}
\end{figure}

\subsection{What is not proved}

We are explicit about the limits of these results, since the value of a
formalization lies in the precision of its claims.

\begin{enumerate}[label=\textup{(\alph*)}]
  \item Theorem~\ref{thm:main} is a rank \emph{lower} bound. Upper bounds
        require descent (Selmer groups, the weak Mordell--Weil theorem),
        which exists in no proof assistant; we make no claim about the exact
        rank of \texttt{389a1}, whose value $2$ is known classically.
  \item Nothing here concerns $L$-functions, analytic ranks, or the
        Birch--Swinnerton-Dyer conjecture. No $L$-function of an elliptic
        curve is defined in our development.
  \item The Mordell--Weil theorem is not used and not proved. All statements
        are about $\operatorname{rank}_{\mathbb{Z}}$ of the point group as an
        abstract $\mathbb{Z}$-module, which is meaningful without finite
        generation.
  \item The classical facts that $\hh$ extends to a positive-definite form on
        $E(\mathbb{Q}) \otimes \mathbb{R}$, and that the regulator of
        \texttt{389a1} is $\approx 0.15246$, are quoted for orientation only
        and are used nowhere in any proof.
\end{enumerate}

\subsection{Related formalization}

\texttt{mathlib}'s elliptic-curve development \cite{mathlib} provides the
Weierstrass type, the group law with a full associativity proof (via the
class group of the coordinate ring), variable change, normal forms and
division polynomials; recent additions include reduction over discrete
valuation rings. Heights appear in \texttt{mathlib} only for number fields
and projective space (\texttt{Mathlib/NumberTheory/Height/}, present on
\texttt{master} after our pin), where the Northcott property is proved and
the point-height and parallelogram law for elliptic curves are marked as
open tasks; M.~Stoll's abstract descent theorem
(\texttt{Mathlib/GroupTheory/Descent.lean}) is likewise a step toward
Mordell--Weil rather than the theorem. Buzzard's FLT project \cite{FLT}
targets a modularity lifting theorem sufficient for Fermat, not the full
modularity theorem. Gross--Zagier, Kolyvagin and Heegner point theory are, to
our knowledge, unformalized anywhere. Our development is independent of all
of these.

\section{Setup}

\subsection{The curve and the na\"ive height}

Throughout, $E$ is the Weierstrass curve
$(a_1,a_2,a_3,a_4,a_6) = (0,1,1,-2,0)$ over $\mathbb{Q}$, i.e.
$y^2 + y = x^3 + x^2 - 2x$, of conductor $389$ and discriminant $389$. In
Lean this is
\begin{center}
\lt{def E389a1 : WeierstrassCurve Q := <0, 1, 1, -2, 0>}
\end{center}
and $E(\mathbb{Q})$ is \lt{E389a1.toAffine.Point}, the inductive type with
constructors \lt{zero} and \lt{some} (the latter carrying a nonsingular
affine point), with \texttt{mathlib}'s \lt{AddCommGroup} instance.

\begin{definition}[\lt{naiveHeight}]
For $q \in \mathbb{Q}$ set $\nh(q) := \max(\lvert \mathrm{num}(q)\rvert,
\mathrm{den}(q)) \in \mathbb{N}$, using Lean's reduced representation. Then
$\nh(q) \ge 1$ always. Define $X \colon E(\mathbb{Q}) \to \mathbb{Q}$ by
$X(O) = 0$ and $X(\text{some at } (x,y)) = x$, and
$h(R) := \log \nh(X(R))$.
\end{definition}

\subsection{The duplication estimate}

\begin{proposition}[\lt{dbl\_x}, \lt{g\_ne\_zero}]\label{prop:dbl}
For every nonsingular rational point $(x,y)$ of $E$ we have
$(2y+1)^2 = g(x)$, and $g$ has no rational zero; consequently no rational
affine point is $2$-torsion, $P + P$ is again affine, and its
$x$-coordinate is $f(x)/g(x)$ with $f,g$ as in~\eqref{eq:fg}.
\end{proposition}

That $g(x) = 4x^3+4x^2-8x+1$ has no rational root is proved by the integer
form of the rational root theorem: a reduced $n/d$ with $4n^3 + 4n^2d -
8nd^2 + d^3 = 0$ forces $n \mid d^3$, hence $n = \pm 1$, and $d \mid 4n^3$,
hence $d \mid 4$; the eight candidates are refuted by \lt{norm\_num}.

\begin{proposition}[\lt{duplication\_height\_bound},
\lt{duplication\_height\_upper}]\label{prop:window}
For every $x \in \mathbb{Q}$,
\[
  \nh(x)^4 \;\le\; 1728\,\nh\!\left(f(x)/g(x)\right), \qquad
  \nh\!\left(f(x)/g(x)\right) \;\le\; 17\,\nh(x)^4 .
\]
\end{proposition}

The upper bound is the triangle inequality on the homogenized numerator and
denominator. The lower bound is the substantive one. Writing $x = a/b$
reduced and
\[
  F(a,b) = a^4 + 4a^2b^2 - 2ab^3 + 3b^4, \qquad
  G(a,b) = 4a^3 + 4a^2b - 8ab^2 + b^3, \qquad D = bG,
\]
so that $x(2R) = F/D$, the extended Euclidean algorithm supplies
\begin{align}
  (-48a^2 - 32ab + 144b^2)F + (12a^3 - 4a^2b + 40ab^2 - 43b^3)G
    &= 389\,b^6, \label{eq:bez1}\\
  (389a^3 + 188a^2b - 520ab^2 + 66b^3)F
    + (-47a^3 - 212a^2b + 108ab^2 - 198b^3)D &= 389\,a^7. \label{eq:bez2}
\end{align}
Both are verified in Lean by \lt{ring}. Since $\gcd(a,b) = 1$ they give
$\gcd(F,D) \mid 389$, and comparing coefficient sums gives
$389\,\nh(x)^7 \le 1728\,\nh(x)^3 \max(\lvert F\rvert, \lvert D\rvert)$,
whence the stated bound after cancelling the content.

\section{The canonical height}

\begin{definition}[\lt{hseq}, \lt{canheight}]
For $R \in E(\mathbb{Q})$ and $n \in \mathbb{N}$ put
$s_n(R) := h(2^n R)/4^n$, and $\hh(R) := \lim_{n} s_n(R)$.
\end{definition}

\begin{proposition}[\lt{hseq\_cauchySeq}, \lt{canheight\_window}]
$\lvert s_{n+1}(R) - s_n(R)\rvert \le \log 1728 / 4^{\,n+1}$, so
$(s_n(R))_n$ is Cauchy and $\hh$ is well defined; moreover
$\lvert \hh(R) - h(R)\rvert \le \tfrac13 \log 1728$.
\end{proposition}

Proposition~\ref{prop:window} gives
$\lvert h(2R) - 4h(R)\rvert \le \log 1728$ for affine $R$ (and trivially for
$R = O$), which after division by $4^{n+1}$ is the displayed step estimate;
Cauchyness follows from the geometric criterion, and the window from the
corresponding tail bound. Applying the window at $2^nR$ and using
$\hh(2^nR) = 4^n \hh(R)$ sharpens it to
\begin{equation}\label{eq:winshift}
  \lvert \hh(R) - s_n(R)\rvert \;\le\; \frac{\log 1728}{3 \cdot 4^n},
\end{equation}
which is the form used for numerics in \S\ref{sec:numerics}.

\begin{figure}[htbp]
\centering
\pffig{fig_canheight_convergence}
\caption{The canonical height on \texttt{389a1}, for the four points that
enter Theorem~\ref{thm:main}. \textbf{(a)} the defining sequence
$s_n(R) = \log\mathcal{H}(x(2^nR))/4^n$; dashed lines mark the limits, and
the shaded funnel is the certified envelope
$\hh(P) \pm 7.625/(3\cdot4^n)$ from~\eqref{eq:winshift}. \textbf{(b)} the
same data as absolute error against that envelope, on a logarithmic scale:
the proved bound is a genuine over-estimate everywhere, by roughly one order
of magnitude, which is what makes the coarse level-$3$ evaluation
of~\S\ref{sec:numerics} sufficient. Computed by exact rational arithmetic
from the mathlib group law; the numerator of $x(2^6P)$ has $582$ decimal
digits.}
\label{fig:conv}
\end{figure}

\begin{theorem}[\lt{canheight\_eq\_zero\_iff\_torsion}]\label{thm:torsion}
$\hh(R) = 0$ if and only if $R$ is a torsion point.
\end{theorem}

\begin{proof}[Proof sketch, as formalized]
If $R$ has finite order its multiples form a finite set, so $h(2^nR)$ is
bounded and $s_n(R) \to 0$. Conversely if $\hh(R) = 0$ then
$\hh(2^nR) = 4^n \hh(R) = 0$ for all $n$, so by the window
$h(2^nR) \le \tfrac13\log 1728$, i.e. $\nh(X(2^nR))^3 \le 1728$ and hence
$\nh(X(2^nR)) \le 12$. A Northcott statement for $\mathbb{Q}$
(\lt{finite\_of\_naiveHeight\_le}: the set of rationals of na\"ive height at
most $B$ is finite, by injecting into a product of finite intervals of
numerators and denominators) together with the fact that each $x$ admits at
most two points (\lt{finite\_equation\_fiber}, from
$(y-y_0)(y+y_0+1) = 0$) shows that $\{2^nR\}_n$ is finite; pigeonholing
gives $2^mR = 2^kR$ with $m < k$, so $(2^k - 2^m)R = 0$ with
$2^k - 2^m \ne 0$.
\end{proof}

\section{The parallelogram law}

\subsection{The addition quadratic}

\begin{proposition}[\lt{xAdd\_sum}, \lt{xAdd\_prod}]\label{prop:addquad}
Let $(x_1,y_1)$, $(x_2,y_2)$ be nonsingular rational points of $E$ with
$x_1 \ne x_2$. Then, with $dd = (x_1-x_2)^2$,
\[
  x(P_1+P_2) + x(P_1-P_2) = \frac{S(x_1,x_2)}{dd}, \qquad
  x(P_1+P_2)\cdot x(P_1-P_2) = \frac{\Pi(x_1,x_2)}{dd},
\]
where
\begin{align*}
  S(x_1,x_2) &= 2x_1^2x_2 + 2x_1x_2^2 + 4x_1x_2 - 4x_1 - 4x_2 + 1,\\
  \Pi(x_1,x_2) &= x_1^2x_2^2 + 4x_1x_2 - x_1 - x_2 + 3 .
\end{align*}
\end{proposition}

The $y$-coordinates cancel entirely. Writing
$N_\pm$ for $dd \cdot x(P_1 \pm P_2)$ as computed from
\texttt{mathlib}'s \lt{slope} and \lt{addX}, and $\mathcal{E}_i$ for the
Weierstrass relations, one has the polynomial identities
\[
  N_+ + N_- - S = 2\mathcal{E}_1 + 2\mathcal{E}_2, \qquad
  N_+ N_- - dd \cdot \Pi = c_1\mathcal{E}_1 + c_2\mathcal{E}_2
\]
with $c_1, c_2$ explicit of degree $3$; both are discharged in Lean by
\lt{linear\_combination}. The homogenizations of $S$ and $\Pi$ are precisely
the bidegree-$(2,2)$ forms whose content and size are controlled as
in~\S2.2, so $x(P_1 \pm P_2)$ are the roots of an integer quadratic
$(DD)T^2 - (Sf)T + (Pf)$ with $\max(\lvert DD\rvert, \lvert Sf\rvert,
\lvert Pf\rvert) \le 17\,\nh(x_1)^2\nh(x_2)^2$.

\subsection{Heights of the roots of an integer quadratic}

\begin{lemma}[\lt{naiveHeight\_mul\_le\_of\_quadratic}]\label{lem:roots}
Let $A,B,C \in \mathbb{Z}$ with $A \ne 0$ and let $u,v \in \mathbb{Q}$
satisfy $A(u+v) = B$ and $Auv = C$. Then
\[
  \nh(u)\,\nh(v) \;\le\; 2\max(\lvert A\rvert, \lvert B\rvert,
  \lvert C\rvert).
\]
\end{lemma}

\begin{proof}[Proof, as formalized]
Write $u = p_1/q_1$, $v = p_2/q_2$ reduced. The rational root theorem applied
to $AT^2 - BT + C$ gives $q_1 \mid A$; writing $A = q_1 m$ and cancelling
$q_1$ in the product and sum relations gives $q_2 \mid m p_1$ and
$q_2 \mid m q_1$, so a B\'ezout combination for $\gcd(p_1,q_1) = 1$ yields
$q_2 \mid m$ and hence $q_1q_2 \mid A$. (This avoids Gauss's lemma and any
prime or three-term gcd reasoning.) Consequently
$(A,B,C) = k\,(q_1q_2,\; p_1q_2 + p_2q_1,\; p_1p_2)$ with $\lvert k\rvert
\ge 1$, so it suffices to bound
$\max(p_1,q_1)\max(p_2,q_2)$, which expands to a maximum of four products,
by twice $M := \max(q_1q_2, \lvert p_1q_2+p_2q_1\rvert, p_1p_2)$. The two
diagonal products are $\le M$ outright. For the cross terms $s = p_1q_2$,
$t = q_1p_2$ one has $st = (p_1p_2)(q_1q_2) \le M^2$, so
$\min(s,t) \le M$, while $\lvert s - t\rvert \le \lvert p_1q_2 +
p_2q_1\rvert \le M$; hence $\max(s,t) \le 2M$.
\end{proof}

\subsection{The quasi-parallelogram law, and the limit}

\begin{proposition}[\lt{point\_two\_sided}, \lt{lognh\_defect\_le}]\label{prop:quasi}
For nonsingular rational points $P,Q$ of $E$ with $x(P) \ne x(Q)$,
\[
  \frac{\nh(x_P)^2\nh(x_Q)^2}{10368} \;\le\;
  \nh\!\left(x(P{+}Q)\right)\nh\!\left(x(P{-}Q)\right) \;\le\;
  34\,\nh(x_P)^2\nh(x_Q)^2,
\]
equivalently
$\lvert h(P{+}Q) + h(P{-}Q) - 2h(P) - 2h(Q)\rvert \le \log 10368$.
\end{proposition}

The upper bound is Proposition~\ref{prop:addquad} with
Lemma~\ref{lem:roots} and the size bound. For the lower bound, apply the
upper bound to the pair $(P{+}Q, P{-}Q)$ --- legitimate because
$x(P{+}Q) \ne x(P{-}Q)$, else the two points would be equal or opposite,
forcing $2Q = O$ or $2P = O$, excluded by Proposition~\ref{prop:dbl} --- and
note that their sum and difference are $2P$ and $2Q$. Combining with the
duplication lower bound of Proposition~\ref{prop:window} gives
$\nh(x_P)^4\nh(x_Q)^4 \le 1728^2 \cdot 34 \cdot
(\nh(x_{P+Q})\nh(x_{P-Q}))^2$, and $1728^2 \cdot 34 = 101\,523\,456 \le
107\,495\,424 = 10368^2$, so extracting square roots in $\mathbb{N}$ gives
the claim.

\begin{theorem}[\lt{canheight\_parallelogram}]\label{thm:par}
If $P, Q, P{+}Q, P{-}Q$ are all non-torsion then
\[
  \hh(P+Q) + \hh(P-Q) \;=\; 2\hh(P) + 2\hh(Q).
\]
\end{theorem}

\begin{proof}[Proof, as formalized]
If $P{\pm}Q$ are non-torsion then $2^nP \ne \pm 2^nQ$ for every $n$ --- else
$2^n(P \mp Q) = O$ --- so $x(2^nP) \ne x(2^nQ)$ and
Proposition~\ref{prop:quasi} applies at $(2^nP, 2^nQ)$. Since
$2^n(P \pm Q) = 2^nP \pm 2^nQ$, dividing by $4^n$ bounds the defect of the
scaled sequences by $\log 10368/4^n \to 0$. The scaled sequences converge to
the four canonical heights, so by uniqueness of limits the defect is $0$.
\end{proof}

\begin{corollary}[\lt{canheight\_zsmul}]\label{cor:mult}
$\hh(kR) = k^2\hh(R)$ for all $k \in \mathbb{Z}$ and non-torsion $R$.
\end{corollary}

\begin{proof}[Proof, as formalized]
Theorem~\ref{thm:par} at $((k{+}2)R, R)$, whose four points are nonzero
multiples of $R$ and hence non-torsion, gives
$\hh((k{+}3)R) = 2\hh((k{+}2)R) + 2\hh(R) - \hh((k{+}1)R)$ for $k \ge 0$.
With base values $\hh(R)$ and $\hh(2R) = 4\hh(R)$ a two-step induction gives
the claim for $k \in \mathbb{N}$; note the recurrence must be indexed from
$1$, since at $k = 0$ the difference point is $O$, which is torsion --- that
step is instead supplied by $\hh(2R) = 4\hh(R)$. Negative $k$ reduces to
positive by $\hh(-R) = \hh(R)$, which holds because negation fixes the
$x$-coordinate, so every term of the defining sequence is unchanged.
\end{proof}

\section{Independence from a nonvanishing regulator}

Define the pairing and the regulator determinant by
\[
  \langle P,Q\rangle := \tfrac12\!\left(\hh(P{+}Q) - \hh(P) - \hh(Q)\right),
  \qquad
  \reg(P,Q) := \hh(P)\hh(Q) - \langle P,Q\rangle^2 .
\]

\begin{proposition}[\lt{regDet\_eq\_zero\_of\_relation}]\label{prop:reldet}
Suppose $P, Q, P{\pm}Q$ are non-torsion and $mP + nQ = O$ with
$m,n \in \mathbb{Z}$, $n \ne 0$. Then $\reg(P,Q) = 0$.
\end{proposition}

\begin{proof}[Proof, as formalized]
From $nQ = -mP$ we get $n(P+Q) = (n-m)P$ and $n(P-Q) = (n+m)P$, so
Corollary~\ref{cor:mult} gives $n^2\hh(P \pm Q) = (n \mp m)^2\hh(P)$.
Substituting into Theorem~\ref{thm:par} --- used exactly once, at the pair
$(P,Q)$ --- yields $m^2\hh(P) = n^2\hh(Q)$; substituting into the definition
of the pairing yields $n\langle P,Q\rangle = -m\hh(P)$. Therefore
\[
  n^2 \reg(P,Q) = \hh(P)\left(n^2\hh(Q)\right) -
  \left(n\langle P,Q\rangle\right)^2 = m^2\hh(P)^2 - m^2\hh(P)^2 = 0,
\]
and $n \ne 0$.
\end{proof}

\begin{remark}
Proposition~\ref{prop:reldet} is where the development departs most sharply
from the textbook treatment. The usual route establishes bi-additivity of
$\langle\cdot,\cdot\rangle$ by polarization, requiring the parallelogram law
at some six pairs of points, each contributing four non-torsion side
conditions in our conditional formulation. Using the hypothesized relation
to reduce everything to multiples of a single point avoids this entirely.
\end{remark}

\begin{figure}[htbp]
\centering
\pffig{fig_regulator_lattice}
\caption{Why $\reg(P,Q) \ne 0$ is the right criterion. Classically $\hh$ is a
positive-definite quadratic form on $E(\mathbb{Q})\otimes\mathbb{R}$, so
$mP+nQ \mapsto m v_1 + n v_2$ realises the subgroup generated by $P$ and $Q$
as a lattice in which $\hh$ is the squared length, and $\reg$ is the squared
area of the fundamental cell. \textbf{(a)} the actual lattice of
\texttt{389a1}: $v_1,v_2$ are reconstructed from the computed
$\hh(P)$, $\hh(Q)$, $\hh(P{+}Q)$ alone, and that $\lVert v_1-v_2\rVert^2$
then agrees with the independently computed $\hh(P{-}Q)$ to
$1.1\times10^{-8}$ is a numerical check on Theorem~\ref{thm:par}.
\textbf{(b)} what a relation does. Holding $v_1$ fixed and sliding $v_2$
towards the line through it raises $\langle P,Q\rangle$ and shrinks the cell
(dashed outlines); a relation such as $2P=3Q$ is the limiting case, in which
$v_2$ is a sub-segment of $v_1$, every $mP+nQ$ lands on one line, and the
area is $0$. Panel~(b) is illustrative and uses no curve data.
\emph{The positive-definiteness
these pictures illustrate is used in no proof}:
Proposition~\ref{prop:reldet} establishes the implication algebraically,
from the parallelogram law and Corollary~\ref{cor:mult} only, and in
particular never needs $\hh \ge 0$.}
\label{fig:lattice}
\end{figure}

\begin{corollary}[\lt{rank\_ge\_two\_of\_regDet\_ne\_zero}]\label{cor:crit}
If $P, Q, P{\pm}Q$ are non-torsion and $\reg(P,Q) \ne 0$ then
$mP + nQ = O$ implies $m = n = 0$, and consequently
$2 \le \operatorname{rank}_{\mathbb{Z}} E(\mathbb{Q})$.
\end{corollary}

For the rank statement, the map
$\mathbb{Z}^2 \to E(\mathbb{Q})$, $(m,n) \mapsto mP + nQ$, is
$\mathbb{Z}$-linear with trivial kernel, and
$\operatorname{rank}_{\mathbb{Z}} \mathbb{Z}^2 = 2$, so
\lt{LinearMap.lift\_rank\_le\_of\_injective} applies.

\section{Certified numerics}\label{sec:numerics}

It remains to verify $\reg(P,Q) > 0$ for $P = (0,0)$, $Q = (1,0)$. Two
devices make this a finite, exact computation.

\emph{The shifted window.} By~\eqref{eq:winshift} at $n = 3$,
$\lvert \hh(R) - s_3(R)\rvert \le \log 1728/192 < 7.625/192 < 0.0398$.

\emph{The dyadic log bracket.} $s_3(R) = \log \nh(x(8R))/64$ where
$\nh(x(8R))$ is an explicit positive integer $\mathcal{H}$. Bracketing
$2^j \le \mathcal{H} < 2^{j+1}$ --- a \lt{norm\_num} fact about integers ---
gives $j\log 2 \le \log \mathcal{H} \le (j+1)\log 2$, and
\texttt{mathlib}'s $\log 2$ bounds \cite{mathlib} pin this to nine decimals.
One bit costs only $\log 2/64 < 0.0109$ at this level. (This avoids
estimating $\log$ of a 26-digit integer directly.)

By Proposition~\ref{prop:dbl}, $x(8R)$ is the third iterate of
$x \mapsto f(x)/g(x)$, computed by \lt{norm\_num}. The values are:

\begin{center}
\begin{tabular}{llrl}
\toprule
point & attained at & $\mathcal{H}\bigl(x(8R)\bigr)$ & bracket \\
\midrule
$P=(0,0)$ & numerator   & $1169154495$                 & $2^{30} \le \mathcal{H} < 2^{31}$\\
$Q=(1,0)$ & numerator   & $11776563836346$             & $2^{43} \le \mathcal{H} < 2^{44}$\\
$P{+}Q$   & numerator   & $30157844295583882647192303$ & $2^{84} \le \mathcal{H} < 2^{85}$\\
$P{-}Q$   & denominator & $7960116832793145801$        & $2^{62} \le \mathcal{H} < 2^{63}$\\
\bottomrule
\end{tabular}
\end{center}

\noindent
The complete fractions --- whose numerators and denominators run to $26$
digits --- are displayed in \lt{RegulatorPositive389a1\_r153.lean}; only the
larger of the two parts, tabulated above, enters $\mathcal{H}$.

Here $x(P{+}Q) = -2$ and $x(P{-}Q) = -1$ are themselves computed from
\texttt{mathlib}'s secant construction, not assumed. The resulting intervals
are
\[
  \hh(P) \in [0.2851, 0.3755], \quad \hh(Q) \in [0.4259, 0.5163],
\]
\[
  \hh(P{+}Q) \in [0.8700, 0.9605], \quad \hh(P{-}Q) \in [0.6317, 0.7221].
\]

All four lower bounds are positive, so by Theorem~\ref{thm:torsion} all four
points are non-torsion --- precisely the side conditions of
Theorem~\ref{thm:par} and Corollary~\ref{cor:crit}, obtained with no separate
argument. Finally $\hh(P)\hh(Q) \ge 0.1214$ while
$\langle P,Q\rangle \in [-0.011, 0.125]$ gives
$\langle P,Q\rangle^2 \le 0.0157$, so
\[
  \reg(P,Q) \;\ge\; 0.1057 \;>\; 0,
\]
and Corollary~\ref{cor:crit} yields Theorem~\ref{thm:main}. \qed

\begin{remark}
The margin is ample: the interval for $\reg$ has width about $0.07$ around a
lower bound of $0.106$, against the classical value $0.15246$. Level $n=3$
was chosen as the smallest sufficient; each further level shrinks both error
sources by a factor of~$4$.
\end{remark}

\begin{figure}[htbp]
\centering
\pffig{fig_certified_windows}
\caption{The four certified enclosures at level $n=3$. For each point the
saturated bar is the dyadic bracket
$[\,j\log 2,\;(j+1)\log 2\,]/4^3$ on $s_3$, obtained from
$2^j \le \mathcal{H}(x(8R)) < 2^{j+1}$ and mathlib's nine-decimal bounds on
$\log 2$; the pale bar adds the analytic window
$\pm\,7.625/(3\cdot 4^3)$ of~\eqref{eq:winshift} to reach $\hh$. The
vertical rule is the true value. Interval arithmetic on the enclosures for
$P$, $Q$, $P{+}Q$ gives
$\langle P,Q\rangle \in [-0.0109, 0.1246]$ and hence
$\reg(P,Q) \ge 0.1057 > 0$, which is all Theorem~\ref{thm:main} requires;
the two error sources are deliberately left coarse, since the conclusion
needs three digits and not thirty.}
\label{fig:windows}
\end{figure}

\section{The rank-one cohort}

The same machinery, instantiated per curve, gives unconditional rank lower
bounds elsewhere. For each curve below the duplication data $(f,g)$, the two
B\'ezout identities, the constant in Proposition~\ref{prop:window}, and a
dyadic chain from the indicated point were computed and verified; the
argument is then Theorem~\ref{thm:torsion} applied to a point whose chain
heights grow, giving a non-torsion point and hence rank $\ge 1$.

\begin{center}
\begin{tabular}{llrr}
\toprule
curve & anchor point & constant & chain height reached\\
\midrule
\texttt{37a1}   & $(0,0)$  & $171$ & $480\,106$\\
\texttt{43a1}   & $(0,0)$  & $139$ & $8\,338\,438$\\
\texttt{53a1}   & $(0,0)$  & $172$ & $369$\\
\texttt{61a1}   & $(1,0)$  & $590$ & $636\,789\,825$\\
\texttt{79a1}   & $(0,0)$  & $298$ & $385$\\
\texttt{83a1}   & $(0,0)$  & $470$ & $83\,281$\\
\texttt{89a1}   & $(0,0)$  & $220$ & $1\,024$\\
\texttt{101a1}  & $(-1,0)$ & $373$ & $28\,981$\\
\texttt{106a1}  & $(2,1)$  & $38\,896$ & $26\,615\,281$\\
\texttt{389a1}  & $(0,0)$  & $1\,728$ & $1\,169\,154\,495$\\
\texttt{5077a1} & $(-2,3)$ & $105\,754$ & $3\,009\,638\,454$\\
\bottomrule
\end{tabular}
\end{center}

The curve \texttt{102a1} is excluded: it has a rational $2$-torsion point
($g(0) = 0$), so the global argument of Proposition~\ref{prop:dbl} fails and
a per-step variant is required. We have not carried this out.

\begin{figure}[htbp]
\centering
\pffig{fig_cohort}
\caption{Scale of the eleven certificates, on a logarithmic axis. Blue is
the na\"ive height reached by the dyadic chain --- the size of the integers
the kernel must actually verify, spanning nearly seven orders of magnitude from
\texttt{53a1} to \texttt{5077a1}. Orange is the duplication constant
$\kappa$ of Proposition~\ref{prop:window}, which controls how long a chain
must run before Theorem~\ref{thm:torsion} applies. The two are largely
independent: \texttt{106a1} needs a large $\kappa$ but a short chain, while
\texttt{61a1} has a small $\kappa$ and a long one. Only \texttt{389a1}
carries the full rank-$2$ argument; the rest stop at rank $\ge 1$.}
\label{fig:cohort}
\end{figure}

\section{Verification data}

All results were checked with Lean~4 toolchain
\texttt{leanprover/lean4:v4.24.0-rc1} against \texttt{mathlib} at commit
\texttt{eed770a4}. At repository HEAD \texttt{0d3b5415}:
\lt{lake build PF} completes in $4658$ jobs and \lt{lake build} in $9301$
jobs, both with zero errors and zero warnings. Every theorem cited in this
paper reports exactly
\[
  [\mathtt{propext},\ \mathtt{Classical.choice},\ \mathtt{Quot.sound}]
\]
under \lt{\#print axioms}; several of the polynomial identities report only
$[\mathtt{propext}, \mathtt{Quot.sound}]$. The development contains no
\lt{sorry}, no \lt{native\_decide} (which would introduce
\lt{Lean.ofReduceBool}), and no project-level \lt{axiom} declarations.

The files, in dependency order, are the following; they are set
ragged-right because the names are unbreakable.

\begin{flushleft}
\lt{MordellWeilRankLowerBound\_r129}, \lt{NaiveHeightQ\_r130},
\lt{DuplicationBezout37a1\_r131}--\lt{E37a1RankOne\_r134} (the rank-one
template), \lt{E43a1RankOne\_r135}--\lt{E106a1RankOne\_r142} (the cohort),
\lt{E389a1RankOne\_r143}, \lt{E5077a1RankOne\_r144},
\lt{DuplicationHeightUpper389a1\_r145}, \lt{CanonicalHeight389a1\_r147},
\lt{QuasiParallelogramCerts389a1\_r148a}--\lt{QuasiParallelogramLower389a1\_r148g},
\lt{PointQuasiParallelogram389a1\_r149},
\lt{CanheightParallelogram389a1\_r150}, \lt{CanheightMultiple389a1\_r151},
\lt{RegulatorIndependence389a1\_r152}, \lt{RegulatorPositive389a1\_r153} and
\lt{E389a1RankTwo\_r154}.
\end{flushleft}

\section{Outlook}

Three extensions appear within reach with the present machinery.

\emph{Rank three.} The curve \texttt{5077a1}, $y^2+y = x^3-7x+6$, is the
smallest-conductor curve of rank $3$ and its duplication data are already
verified (\lt{E5077a1RankOne\_r144}). Extending
Proposition~\ref{prop:reldet} from two points to three is not immediate: the
trick of reducing a relation to multiples of a single point does not
generalize, so the $3 \times 3$ case appears to require genuine
bi-additivity of the pairing, i.e. the polarization argument this development
was able to avoid.

\emph{Torsion.} With $\hh$ available, $\hh(R) = 0$ is a computable
obstruction; combined with a bound on the possible torsion orders it would
give exact torsion subgroups for individual curves.

\emph{Upstreaming.} The general components --- Northcott for $\mathbb{Q}$,
the height of the roots of an integer quadratic (Lemma~\ref{lem:roots}), and
the abstract shape ``two-sided duplication estimate $\Rightarrow$ canonical
height'' --- are curve-independent and would fit \texttt{mathlib}'s
\texttt{NumberTheory/Height} directory, where the elliptic point height is
currently an open task.

\section*{Acknowledgements}

The formalization was carried out in collaboration with Claude (Anthropic),
which wrote and verified the Lean sources under the author's direction; all
mathematical claims were independently rebuilt and axiom-audited before being
recorded. The author thanks the \texttt{mathlib} community for the
elliptic-curve group law, without which none of this would have been
possible.

\begin{thebibliography}{9}

\bibitem{FLT}
K.~Buzzard et al.,
\emph{The Fermat's Last Theorem project},
\url{https://github.com/ImperialCollegeLondon/FLT}, 2024--.

\bibitem{cremona}
J.~E.~Cremona,
\emph{Algorithms for Modular Elliptic Curves},
2nd ed., Cambridge University Press, 1997.

\bibitem{lmfdb}
The LMFDB Collaboration,
\emph{The $L$-functions and Modular Forms Database},
\url{https://www.lmfdb.org}, 2026.

\bibitem{mathlib}
The mathlib Community,
\emph{The Lean mathematical library},
CPP 2020, 367--381; repository \url{https://github.com/leanprover-community/mathlib4}.

\bibitem{silverman}
J.~H.~Silverman,
\emph{The Arithmetic of Elliptic Curves},
2nd ed., Graduate Texts in Mathematics 106, Springer, 2009.

\bibitem{tate}
J.~Tate,
\emph{Variation of the canonical height of a point depending on a parameter},
Amer. J. Math. \textbf{105} (1983), 287--294.

\end{thebibliography}

\end{document}
